% =====================================================================
%  Planning projet — Groupe 69 « Conception d'une Électronique Modulaire (Bus CAN) »
%  ENSEA 2A — MANAGEMENT_PROJET_S7 — 2026-2027
%  Version 3.0 — 22 septembre 2026
%  Compiler avec pdflatex.
% =====================================================================
\documentclass[10pt,a4paper,landscape]{article}

\usepackage[utf8]{inputenc}
\usepackage[T1]{fontenc}
\usepackage[french]{babel}
\usepackage{lmodern}
\usepackage[margin=1.2cm]{geometry}
\usepackage{pgfgantt}
\usepackage{booktabs}
\usepackage{tabularx}
\usepackage{array}
\usepackage{xcolor}
\usepackage{fancyhdr}
\usepackage{colortbl}
\usepackage{amssymb}

\definecolor{canblue}{HTML}{1F3864}
\definecolor{canamber}{HTML}{FFB627}
\definecolor{greyrow}{HTML}{F2F2F2}

% ---------- Une couleur par personne ----------
\definecolor{cHugo}{HTML}{2E5E8C}
\definecolor{cDylan}{HTML}{C2661F}
\definecolor{cSamuel}{HTML}{3F8F6B}
\definecolor{cTous}{HTML}{808080}

\newcommand{\bH}[3]{\ganttbar[bar/.append style={fill=cHugo}]{#1}{#2}{#3}}
\newcommand{\bD}[3]{\ganttbar[bar/.append style={fill=cDylan}]{#1}{#2}{#3}}
\newcommand{\bS}[3]{\ganttbar[bar/.append style={fill=cSamuel}]{#1}{#2}{#3}}
\newcommand{\bT}[3]{\ganttbar[bar/.append style={fill=cTous}]{#1}{#2}{#3}}

\pagestyle{fancy}
\fancyhf{}
\lhead{\small Planning projet — Groupe 69 — Bus CAN}
\rhead{\small v3.0 — 22 septembre 2026}
\cfoot{\small \thepage}
\renewcommand{\headrulewidth}{0.4pt}

\newcolumntype{L}[1]{>{\raggedright\arraybackslash}p{#1}}
\newcolumntype{C}[1]{>{\centering\arraybackslash}p{#1}}

\begin{document}

% =====================================================================
\begin{center}
  {\color{canblue}\Large\bfseries Planning prévisionnel du projet\par}
  \vspace{1pt}
  {\small Groupe 69 — Conception d'une Électronique Modulaire (Bus CAN) — Samuel Gomri · Dylan Kiati · Hugo Mathoux\par}
  {\color{canamber}\rule{0.4\textwidth}{1.2pt}\par}
\end{center}

\vspace{-10pt}
\noindent\scriptsize
Séances le mardi après-midi (4 h) — semaine 1 = semaine du 22 septembre 2026 (séance n\textsuperscript{o}2), rendu final le 10 janvier 2027.
Semaines 6, 14 et 15 supposées non travaillées : \textbf{à confirmer}.
\normalsize

\vspace{1pt}

% =====================================================================
%  GANTT — PARTIE 1 : pilotage, ICD, banc de test
% =====================================================================
\begin{center}
\begin{ganttchart}[
    hgrid, vgrid,
    x unit=1.22cm,
    y unit title=0.48cm,
    y unit chart=0.325cm,
    title height=1,
    title label font=\scriptsize\bfseries,
    bar label font=\scriptsize,
    group label font=\scriptsize\bfseries,
    milestone label font=\scriptsize\bfseries,
    bar height=0.56,
    group/.append style={draw=none, fill=black!65},
    group left shift=0, group right shift=0, group top shift=0.38, group height=0.16,
    bar/.append style={draw=none},
    milestone/.append style={fill=canamber, draw=none},
    bar label node/.append style={text width=5.8cm, align=right},
    group label node/.append style={text width=5.8cm, align=right},
    milestone label node/.append style={text width=5.8cm, align=right},
  ]{1}{16}

  \gantttitle{Sept.}{1}
  \gantttitle{Octobre}{4}
  \gantttitle{Novembre}{4}
  \gantttitle{Décembre}{4}
  \gantttitle{Janvier}{3} \\
  \gantttitlelist{1,...,16}{1} \\

  % ---------- LOT 0 : pilotage ----------
  \ganttgroup{L0 — Pilotage projet}{1}{16} \\
  \bH{Dépôt partagé : Git + Drive}{1}{1} \\
  \bH{Matrice RACI}{1}{2} \\
  \bH{Registre de risques}{2}{3} \\
  \bD{Nomenclature et suivi des coûts}{1}{4} \\
  \bT{Compte rendu de séance (tournant)}{1}{16} \\
  \bH{Mise à jour du planning}{1}{16} \\
  \bT{Revues d'avancement (encadrants)}{5}{5} \bT{}{9}{9} \bT{}{13}{13} \\
  \bS{Carnet de bord DD et ACV}{9}{12} \\

  % ---------- LOT 1 : standard et ICD ----------
  \ganttgroup{L1 — Standard \& ICD}{1}{13} \\
  \bS{Diffusion des ICD vierges}{1}{1} \\
  \bS{Relances et collecte des retours}{2}{9} \\
  \bH{ICD Ecomarathon v1}{4}{12} \\
  \bH{ICD Gyropode v1}{4}{12} \\
  \bD{ICD Formula Student v1}{4}{12} \\
  \bS{ICD COBOT v1}{4}{12} \\
  \bS{ICD MIR1000 v1}{4}{12} \\
  \bT{Accompagnement des groupes prêts à intégrer}{11}{13} \\

  % ---------- LOT 2 : banc de test ----------
  \ganttgroup{L2 — Banc de test (matériel)}{1}{11} \\
  \bD{Inventaire et réservation du matériel}{1}{1} \\
  \bD{Commande du complément}{2}{4} \\
  \bD{Alimentation et masse commune}{2}{3} \\
  \bD{Câblage bus, bornier, terminaisons}{2}{4} \\
  \bD{Points de test oscilloscope}{4}{5} \\
  \bD{Schéma électrique du banc}{5}{7} \\
  \bS{Plaque de montage}{4}{7} \\
  \bH{Nœud 1 : STM32 + capteurs}{3}{5} \\
  \bH{Nœud 2 : STM32 récepteur}{3}{5} \\
  \bD{Nœud 3 : ESP32}{5}{8} \\
  \bS{Nœud 4 : Arduino + MCP2515}{7}{10} \\
  \bD{Nœud défaillant (injection de fautes)}{9}{11} \\
  \bS{IHM, écran de contrôle}{7}{11} \\
  \bS{Analyseur USB-CAN + SavvyCAN}{5}{8}

\end{ganttchart}
\end{center}

\vspace{2pt}
\noindent\footnotesize
\textcolor{cHugo}{$\blacksquare$} Hugo \quad
\textcolor{cDylan}{$\blacksquare$} Dylan \quad
\textcolor{cSamuel}{$\blacksquare$} Samuel \quad
\textcolor{cTous}{$\blacksquare$} Ensemble \quad
\textcolor{canamber}{$\blacklozenge$} Jalon
\normalsize

\clearpage

% =====================================================================
%  GANTT — PARTIE 2 : logiciel, mesures, livrables
% =====================================================================
\begin{center}
\begin{ganttchart}[
    hgrid, vgrid,
    x unit=1.22cm,
    y unit title=0.48cm,
    y unit chart=0.325cm,
    title height=1,
    title label font=\scriptsize\bfseries,
    bar label font=\scriptsize,
    group label font=\scriptsize\bfseries,
    milestone label font=\scriptsize\bfseries,
    bar height=0.56,
    group/.append style={draw=none, fill=black!65},
    group left shift=0, group right shift=0, group top shift=0.38, group height=0.16,
    bar/.append style={draw=none},
    milestone/.append style={fill=canamber, draw=none},
    bar label node/.append style={text width=5.8cm, align=right},
    group label node/.append style={text width=5.8cm, align=right},
    milestone label node/.append style={text width=5.8cm, align=right},
  ]{1}{16}

  \gantttitle{Sept.}{1}
  \gantttitle{Octobre}{4}
  \gantttitle{Novembre}{4}
  \gantttitle{Décembre}{4}
  \gantttitle{Janvier}{3} \\
  \gantttitlelist{1,...,16}{1} \\

  % ---------- LOT 3 : logiciel ----------
  \ganttgroup{L3 — Briques logicielles}{2}{13} \\
  \bH{Calcul du bit timing 500 kbit/s}{2}{4} \\
  \bH{Réglage du point d'échantillonnage}{3}{5} \\
  \bH{Filtres de réception}{5}{7} \\
  \bH{Bibliothèque encodage / décodage}{3}{8} \\
  \bH{Tests unitaires de la bibliothèque (PC)}{6}{8} \\
  \bD{Portage ESP32 et Arduino}{7}{10} \\
  \bH{Générateur de code depuis l'ICD}{7}{11} \\
  \bH{Export DBC}{9}{11} \\
  \bS{Notice d'utilisation}{11}{13} \\

  % ---------- LOT 4 : mesures et validation ----------
  \ganttgroup{L4 — Mesures \& validation}{6}{13} \\
  \bH{Cahier de recette (procédure de test)}{6}{8} \\
  \bD{Oscillo : signal différentiel}{8}{10} \\
  \bD{Effet d'une terminaison manquante}{8}{10} \\
  \bD{Arbitrage à l'analyseur logique}{9}{11} \\
  \bD{Analyse sécurité et réglementaire du banc}{9}{10} \\
  \bD{Injection de fautes, bus-off}{10}{12} \\
  \bH{Mesure de charge du bus}{11}{13} \\
  \bT{Exécution de la recette}{12}{13} \\
  \bS{Dossier de mesures}{10}{13} \\

  % ---------- LOT 5 : livrables évalués ----------
  \ganttgroup{L5 — Livrables évalués}{10}{16} \\
  \bT{Rapport écrit}{10}{16} \\
  \bS{Support oral}{12}{16} \\
  \bT{Répétition de la soutenance}{16}{16} \\
  \bT{Bilan de clôture, écart prévu / réalisé}{16}{16} \\

  % ---------- JALONS ----------
  \ganttmilestone{J1 — Matériel réservé, dépôt en place}{1} \\
  \ganttmilestone{J2 — 1\textsuperscript{re} trame échangée}{5} \\
  \ganttmilestone{J3 — 3 familles de µC dialoguent}{8} \\
  \ganttmilestone{J4 — Banc complet + IHM}{11} \\
  \ganttmilestone{J5 — Recette passée, 5 ICD v1}{13} \\
  \ganttmilestone{J6 — Rendu final}{16}

\end{ganttchart}
\end{center}

\vspace{2pt}
\noindent\footnotesize
\textcolor{cHugo}{$\blacksquare$} Hugo \quad
\textcolor{cDylan}{$\blacksquare$} Dylan \quad
\textcolor{cSamuel}{$\blacksquare$} Samuel \quad
\textcolor{cTous}{$\blacksquare$} Ensemble \quad
\textcolor{canamber}{$\blacklozenge$} Jalon
\normalsize

\clearpage

% =====================================================================
%  CE QUE CONTIENT LE BANC
% =====================================================================
\section*{Composition du banc de test}

\noindent\small
Quatre nœuds de trois familles de microcontrôleurs différentes sur un même bus : c'est
l'hétérogénéité qui démontre que le standard est réutilisable. Chaque nœud est une brique
que les autres groupes peuvent reprendre telle quelle.
\normalsize

\vspace{6pt}
\noindent
\begin{tabularx}{\textwidth}{L{2.0cm} L{3.4cm} L{3.4cm} L{3.0cm} X}
  \toprule
  \textbf{Nœud} & \textbf{Microcontrôleur} & \textbf{Transceiver} & \textbf{Rôle} & \textbf{Ce qu'il démontre} \\
  \midrule
  \rowcolor{greyrow}
  1 & STM32 Nucleo (CAN intégré) & MCP2562 (5 V, VIO 3,3 V) & Émetteur : potentiomètre (vitesse) + thermistance (température) & Le cas typique des groupes partenaires \\
  2 & STM32 Nucleo & MCP2562 & Récepteur + relais vers l'IHM & Réception filtrée, décodage \\
  \rowcolor{greyrow}
  3 & ESP32 (CAN intégré, TWAI) & SN65HVD230 (3,3 V natif) & Émetteur / récepteur & Interopérabilité 3,3 V $\leftrightarrow$ 5 V sans adaptateur \\
  4 & Arduino Uno (pas de CAN intégré) & Module MCP2515 (SPI) & Émetteur basique & Que le standard marche aussi avec un contrôleur externe à 5 euros \\
  \rowcolor{greyrow}
  — & PC + CANable & — & Analyseur (SavvyCAN) & Observation et mesure du trafic réel \\
  \bottomrule
\end{tabularx}

\vspace{10pt}

\noindent
\begin{tabularx}{\textwidth}{L{5.0cm} X}
  \toprule
  \textbf{Infrastructure du banc} & \textbf{Pourquoi c'est indispensable} \\
  \midrule
  \rowcolor{greyrow}
  Masse commune à tous les nœuds &
  Le CAN est différentiel mais tolère au maximum $\pm$12 V de mode commun. Sans masse reliée entre les cartes, la communication échoue de façon erratique. Cause d'échec n\textsuperscript{o}1 sur un banc improvisé. \\
  Alimentation unique 5 V / 3,3 V &
  Évite d'alimenter chaque carte par un PC différent, ce qui créerait justement des masses flottantes. \\
  \rowcolor{greyrow}
  Terminaisons 120\,$\Omega$ sur jumper &
  Permet de les retirer volontairement pour filmer les réflexions à l'oscilloscope (mesure L4). Une terminaison soudée en dur interdit cette démonstration. \\
  Bornier de dérivation &
  Brancher ou débrancher un nœud en 10 secondes, sans redémonter le bus — nécessaire pour l'injection de fautes. \\
  \rowcolor{greyrow}
  Points de test CAN\_H / CAN\_L &
  Sonde d'oscilloscope accessible sans débrancher quoi que ce soit. \\
  Plaque de montage &
  Le banc survit d'une séance à l'autre et se transporte. Sans support, tout est à recâbler chaque mardi. \\
  \rowcolor{greyrow}
  Schéma électrique et nomenclature &
  Sans eux, le banc n'est pas reproductible par un autre groupe — et ce n'est donc pas une brique réutilisable. \\
  \bottomrule
\end{tabularx}

\clearpage

% =====================================================================
%  CALENDRIER
% =====================================================================
\section*{Calendrier des séances}

\noindent
\begin{tabularx}{\textwidth}{C{1.5cm} C{2.4cm} C{1.6cm} X}
  \toprule
  \textbf{Semaine} & \textbf{Date (mardi)} & \textbf{Séance} & \textbf{Objectif principal} \\
  \midrule
  \rowcolor{greyrow} 1 & 22 sept. 2026 & 2 & \textbf{J1} — réserver le matériel, vérifier les transceivers ; créer le dépôt partagé ; diffuser les ICD \\
  2 & 29 sept. 2026 & 3 & Alimentation, masse commune, câblage du bus ; calcul du bit timing ; nomenclature \\
  \rowcolor{greyrow} 3 & 6 oct. 2026 & 4 & Firmware des nœuds 1 et 2 ; démarrage de la bibliothèque \\
  4 & 13 oct. 2026 & 5 & Points de test ; réglage du point d'échantillonnage ; plaque de montage \\
  \rowcolor{greyrow} 5 & 20 oct. 2026 & 6 & \textbf{J2} — une trame circule entre 2 STM32 et est décodée ; 1\textsuperscript{re} revue encadrants \\
  6 & 27 oct. 2026 & — & \textit{Vacances (à confirmer)} \\
  \rowcolor{greyrow} 7 & 3 nov. 2026 & 7 & Intégrer l'ESP32 ; analyseur USB-CAN ; schéma électrique ; démarrer l'IHM \\
  8 & 10 nov. 2026 & 8 & \textbf{J3} — 3 familles de µC dialoguent ; tests unitaires passés ; premières mesures oscillo \\
  \rowcolor{greyrow} 9 & 17 nov. 2026 & 9 & Nœud Arduino ; export DBC ; mesures d'arbitrage ; analyse sécurité ; 2\textsuperscript{e} revue \\
  10 & 24 nov. 2026 & 10 & Nœud défaillant ; injection de fautes ; démarrage du rapport \\
  \rowcolor{greyrow} 11 & 1\textsuperscript{er} déc. 2026 & 11 & \textbf{J4} — banc complet 4 nœuds + IHM ; accompagnement des groupes prêts \\
  12 & 8 déc. 2026 & 12 & Consolidation des 5 ICD ; mesure de charge ; exécution de la recette \\
  \rowcolor{greyrow} 13 & 15 déc. 2026 & 13 & \textbf{J5} — recette passée, 5 ICD en v1 ; 3\textsuperscript{e} revue encadrants \\
  14--15 & 22--29 déc. 2026 & — & \textit{Vacances} — rédaction du rapport à distance \\
  \rowcolor{greyrow} 16 & 5 janv. 2027 & 14 & \textbf{J6} — finalisation du rapport, répétition de l'oral, bilan de clôture \\
  \multicolumn{4}{@{}l}{\textbf{Rendu final : 10 janvier 2027}} \\
  \bottomrule
\end{tabularx}

\vspace{10pt}

% =====================================================================
%  RISQUES
% =====================================================================
\section*{Risques sur le planning}

\noindent
\begin{tabularx}{\textwidth}{L{5.4cm} C{1.5cm} C{1.5cm} X}
  \toprule
  \textbf{Risque} & \textbf{Probab.} & \textbf{Impact} & \textbf{Parade intégrée au planning} \\
  \midrule
  \rowcolor{greyrow}
  Matériel pris par un autre groupe & Élevée & Fort &
  Réservation en semaine 1, avant toute autre tâche. Jalon J1 dès la séance 2. \\
  Transceivers inadaptés (5 V sans VIO) & Moyenne & Moyen &
  Vérification des références à l'inventaire ; commande du complément semaines 2 à 4, avant J2. \\
  \rowcolor{greyrow}
  Bit timing incohérent entre familles de µC & Moyenne & Fort &
  Tâche isolée dès la semaine 2, avant tout firmware. Point d'échantillonnage à 87,5 \% imposé à tous les nœuds. \\
  Masses non reliées entre les cartes & Élevée & Fort &
  Alimentation unique et masse commune en semaine 2, avant le premier essai de communication. \\
  \rowcolor{greyrow}
  Retard des groupes partenaires sur les ICD & Élevée & Faible &
  Le chemin critique passe par L2 et L3, qui ne dépendent d'aucun groupe externe. \\
  Concentration des tâches logicielles critiques sur une seule personne & Moyenne & Fort &
  Portage ESP32/Arduino et notice délégués ; la bibliothèque et ses tests unitaires peuvent avancer à distance, sans matériel. \\
  \rowcolor{greyrow}
  Sous-estimation du développement logiciel & Moyenne & Fort &
  Bibliothèque démarrée en semaine 3, en parallèle du banc, avec 6 semaines de marge avant J3. \\
  Dérive du périmètre (demandes des groupes hors standard) & Moyenne & Moyen &
  Tout ajout passe par une nouvelle version de l'ICD concerné et une ligne dans le registre de risques ; rien n'est accepté oralement. \\
  \bottomrule
\end{tabularx}

\vspace{10pt}

% =====================================================================
%  SUIVI
% =====================================================================
\section*{Comment tenir ce planning à jour}

\noindent
\begin{minipage}{\textwidth}
\small
\textbf{En fin de chaque séance}, cinq minutes suffisent :

\begin{enumerate}\itemsep2pt
  \item Noter ce qui a réellement été fait, et ce qui a glissé.
  \item Corriger les dates dans ce fichier en \textbf{conservant l'ancienne version} (v3.0, v3.1\ldots) — c'est l'écart entre le prévu et le réalisé qui démontre la maîtrise du projet.
  \item Reporter tout nouveau risque dans le tableau ci-dessus.
  \item Mettre à jour le journal de l'ICD du groupe rencontré.
\end{enumerate}

\medskip
\textbf{Chemin critique} : réservation du matériel $\rightarrow$ masse et câblage $\rightarrow$ bit timing $\rightarrow$ firmware nœuds 1-2 $\rightarrow$ bibliothèque $\rightarrow$ banc complet $\rightarrow$ recette $\rightarrow$ rapport.
Tout retard sur cette chaîne décale la fin. Les ICD (L1) ont volontairement de la marge : ils ne doivent jamais bloquer le reste.
\normalsize
\end{minipage}

\end{document}